\begin{progprob}
    % No submission.
    % \iffalse
    In lecture, we saw that Kruskal's algorithm can be used to cluster a weighted graph. The name for this approach is
    \emph{single linkage clustering}.

    In a file named \python{slc.py}, write a function \python{slc(graph, d, k)}
    which accepts the following arguments:
    \begin{itemize}
        \item \python{graph}: An instance of \python{dsc40graph.UndirectedGraph}.
        \item \python{d}: A function of one argument: a tuple containing a pair of nodes. It returns
            the distance (or dissimilarity) between them. See the example below.
        \item \python{k}: A positive integer describing the number of clusters which should
            be found.
    \end{itemize}
    The function should perform single linkage clustering using Kruskal's algorithm
    and it should return a \python{frozenset} of $k$ \python{frozenset}s, each
    representing a cluster of the graph.

    Example:

    \begin{minted}[autogobble]{python}
        >>> g = dsc40graph.UndirectedGraph()
        >>> edges = [('a', 'b'), ('a', 'c'), ('c', 'd'), ('b', 'd')]
        >>> for edge in edges: g.add_edge(*edge)
        >>> def d(edge):
        ...     u, v = sorted(edge)
        ...     return {
        ...         ('a', 'b'): 1,
        ...         ('a', 'c'): 4,
        ...         ('b', 'd'): 3,
        ...         ('c', 'd'): 2,
        ...     }[(u, v)]
        >>> slc(g, d, 2)
        frozenset({frozenset({'a', 'b'}), frozenset({'c', 'd'})})
    \end{minted}

    Note: to implement Kruskal's algorithm, you'll need an implementation of a Disjoint
    Set Forest data structure. We've uploaded a simple one here:
    \url{https://gist.github.com/eldridgejm/983d6ce03a82bf295599e9880ef02bab}

    You can copy and paste this into \python{slc.py}, or put it in a separate file that
    is imported; if you do this, make sure to upload that file alongside \python{slc.py}.

    \begin{soln}
    	\inputminted{python}{\thisdir/include-solution/slc_sol.py}
    \end{soln}
    % \fi
\end{progprob}
